% ---------------------------------------------------------------------
% 2.13 Modelado del lenguaje y predicción de palabras
% ---------------------------------------------------------------------
\section[MODELADO DEL LENGUAJE Y PREDICCIÓN DE PALABRAS]{Modelado del lenguaje y predicción de
palabras}
\label{sec:lenguaje}

\subsection{Predicción de palabras en la comunicación asistida}

Aun con una detección perfecta de la P300, escribir carácter por carácter en un P300 Speller es un
proceso lento, pues cada selección requiere varios segundos. La predicción de palabras, ampliamente
utilizada en los sistemas de comunicación aumentativa, reduce el número de selecciones necesarias
al ofrecer al usuario una lista de palabras probables a partir de los caracteres que ya ha escrito,
de modo que puede completar una palabra con una sola selección adicional. Ryan et al. integraron
este mecanismo en un P300 Speller agregando a la matriz una región con las palabras sugeridas, y
encontraron que los participantes completaban las frases en menos tiempo que con el deletreador
convencional, aunque con una exactitud por selección ligeramente menor \cite{ryan2011}. El
beneficio de estos sistemas suele medirse mediante el ahorro de selecciones,
\begin{equation}
\text{AS} = \frac{S_{\text{sin}} - S_{\text{con}}}{S_{\text{sin}}}\times100\%,
\label{eq:ahorro}
\end{equation}
donde $S_{\text{sin}}$ y $S_{\text{con}}$ son el número de selecciones necesarias para escribir un
mismo texto sin y con predicción, respectivamente.

\subsection{Modelos de lenguaje probabilísticos}

Un modelo de lenguaje asigna una probabilidad a cada secuencia de palabras $w_1, \dots, w_m$, lo
que permite estimar cuál es la palabra, o el carácter, más probable dado lo que se ha escrito. Por
la regla de la cadena, esta probabilidad puede descomponerse como el producto de probabilidades
condicionales, y los modelos de $n$-gramas la aproximan suponiendo que cada palabra depende
únicamente de las $n-1$ anteriores, supuesto conocido como de Markov \cite{jurafsky2025}:
\begin{equation}
P(w_1,\dots,w_m) = \prod_{i=1}^{m} P(w_i \mid w_1,\dots,w_{i-1}) \approx \prod_{i=1}^{m} P(w_i \mid w_{i-n+1},\dots,w_{i-1}).
\label{eq:ngramas}
\end{equation}
Las probabilidades condicionales se estiman contando las frecuencias de aparición en un corpus de
texto, y dado que muchas secuencias válidas no aparecen en él, se aplican técnicas de suavizado que
reservan parte de la probabilidad para eventos no observados \cite{jurafsky2025}. La calidad de un
modelo de lenguaje se evalúa comúnmente mediante la perplejidad, $PP = P(w_1,\dots,w_m)^{-1/m}$,
que puede interpretarse como el número promedio de opciones entre las que el modelo duda en cada
paso, concepto que se relaciona directamente con la noción de entropía de la teoría de la
información de Shannon \cite{shannon1948,jurafsky2025}.

\subsection{Modelos de lenguaje neuronales}

Los modelos de $n$-gramas están limitados a un contexto corto y fijo, por lo que fueron
reemplazados progresivamente por modelos neuronales. Las redes recurrentes, y en particular las
LSTM, permitieron condicionar la predicción en un contexto de longitud variable
\cite{hochreiter1997,jurafsky2025}, y posteriormente la arquitectura Transformer, basada
exclusivamente en mecanismos de atención, eliminó la recurrencia y permitió entrenar modelos
mucho más grandes de forma paralela \cite{vaswani2017}. Sobre esta arquitectura se construyen los
modelos de lenguaje de gran escala (LLM), entrenados con enormes cantidades de texto, que son
capaces de realizar tareas nuevas a partir de unas cuantas instrucciones o ejemplos
\cite{brown2020}.

\subsection{Integración del lenguaje en el P300 Speller}

La información del lenguaje puede incorporarse al P300 Speller de dos maneras: como sugerencias
explícitas que el usuario selecciona, o como una probabilidad previa que se combina con los
puntajes del clasificador para decidir el carácter. Speier et al. siguieron la segunda vía,
integrando un modelo de lenguaje en el proceso de clasificación dinámica, y reportaron mejoras
tanto en la exactitud como en la tasa de transferencia de información \cite{speier2012}. Más
recientemente, Hong et al. propusieron ChatBCI, un P300 Speller que utiliza un LLM para predecir
palabras según el contexto de la frase, logrando una reducción del tiempo promedio de escritura de
62.14\% \cite{hong2026}. La Tabla~\ref{tab:lenguaje} compara estos enfoques.

\begin{table}[htbp]
\centering
\caption{Trabajos que integran modelos de lenguaje en el P300 Speller}
\label{tab:lenguaje}
\small
\begin{tabularx}{\textwidth}{L{3.0cm}L{3.2cm}YY}
\toprule
\textbf{Trabajo} & \textbf{Modelo de lenguaje} & \textbf{Forma de integración} & \textbf{Resultado} \\
\midrule
Ryan et al., 2011 \cite{ryan2011} & Predicción de palabras por diccionario & Región adicional con palabras sugeridas en la interfaz. & Menor tiempo para completar frases, con exactitud ligeramente menor. \\
Speier et al., 2012 \cite{speier2012} & Modelo probabilístico de lenguaje natural & Probabilidad previa combinada con la clasificación dinámica. & Mejor exactitud y tasa de transferencia. \\
Hong et al., 2026 \cite{hong2026} & LLM & Predicción contextual de palabras. & Reducción del tiempo promedio de 62.14\%. \\
\bottomrule
\end{tabularx}
\par\smallskip{\footnotesize Elaboración propia.}
\end{table}
